%%% FIELD cited-strassen-statement
Let \(U\) and \(V\) be finite sets, let \(R\subseteq U\times V\), and let \(p\) and \(q\) be probability measures on \(U\) and \(V\).  For every \(\varepsilon\geq0\), there is a coupling \(\gamma\) of \(p\) and \(q\) satisfying \(\gamma(R)\geq1-\varepsilon\) if and only if
\[
p(C)\leq q(N_R(C))+\varepsilon
\quad\text{for every }C\subseteq U,
\]
where \(N_R(C)=\{v\in V:\exists u\in C\text{ such that }(u,v)\in R\}\).  This equivalence permits zero atoms and arbitrary real-valued probability masses.
%%% FIELD hall-replacement-statement
For every \(S\subseteq E\),
\[
h_{\mu}(S)=\max_{A\subseteq X}\{\mu_1(A)-\mu_0(N_{S^{\mathsf c}}(A))\}=\max\!\left\{0,\max_{\substack{\varnothing\neq A,B\subseteq X\\A-B\subseteq S}}\bigl(\mu_1(A)+\mu_0(B)-1\bigr)\right\}.
\]
Both maxima and the minimum defining \(h_{\mu}(S)\) are attained, including when marginal cells are zero.
%%% FIELD hall-proof
Put \(p=\mu_1\), \(q=\mu_0\), and define the outside relation
\[
R_S=\{(i,j)\in X\times X:i-j\notin S\}.
\]
For \(A\subseteq X\), its neighborhood under \(R_S\) is exactly \(N_{S^{\mathsf c}}(A)\) by \(\text{def:complement-neighborhood}\).  Define
\[
\delta_S=\max_{A\subseteq X}\{p(A)-q(N_{S^{\mathsf c}}(A))\}.
\]
The maximum exists because \(X\) is finite, and \(\delta_S\geq0\) because the term for \(A=\varnothing\) is zero.  Apply \(\text{lem:finite-strassen-deficiency}\) to \(R_S\).  If \(\gamma\in\Gamma(p,q)\), then \(\gamma(R_S)=1-\gamma\{(i,j):i-j\in S\}\).  The necessity direction of that lemma, with \(\varepsilon=\gamma\{i-j\in S\}\), gives
\[
p(A)-q(N_{S^{\mathsf c}}(A))\leq \gamma\{i-j\in S\}
\quad\text{for every }A\subseteq X.
\]
Consequently \(\delta_S\leq h_{\mu}(S)\).  Conversely, the defining inequalities for \(\delta_S\) satisfy
\[
p(A)\leq q(N_{S^{\mathsf c}}(A))+\delta_S
\quad\text{for every }A\subseteq X.
\]
The sufficiency direction of the cited lemma therefore supplies a coupling \(\gamma^*\in\Gamma(p,q)\) with \(\gamma^*(R_S)\geq1-\delta_S\), hence with selected mass at most \(\delta_S\).  The preceding lower bound forces equality, so \(h_{\mu}(S)=\delta_S\), and \(\gamma^*\) attains the minimum.  This argument uses neither division by an atom nor positivity of any cell.

It remains to prove the rectangle form.  For any \(A\subseteq X\), set
\[
B_A=X\setminus N_{S^{\mathsf c}}(A).
\]
If \(i\in A\) and \(j\in B_A\), then \(i-j\in S\); therefore \(A-B_A\subseteq S\), and
\[
p(A)-q(N_{S^{\mathsf c}}(A))=p(A)+q(B_A)-1.
\]
If the left side is positive, then both \(A\) and \(B_A\) are nonempty.  If it is nonpositive, it is dominated by the explicit zero in the claimed maximum.  Conversely, whenever nonempty \(A,B\subseteq X\) satisfy \(A-B\subseteq S\), no point of \(B\) belongs to \(N_{S^{\mathsf c}}(A)\).  Hence \(B\subseteq B_A\), and nonnegativity of \(q\) yields
\[
p(A)+q(B)-1\leq p(A)+q(B_A)-1
=p(A)-q(N_{S^{\mathsf c}}(A))\leq\delta_S.
\]
The two finite maxima are thus equal.  Finally, \(\Gamma(p,q)\) is nonempty because it contains the product table \(p_iq_j\), and it is a closed bounded polytope in \(\mathbb R^{L^2}\); this also independently verifies compact attainment of the continuous linear objective, including at boundary marginal laws.
%%% FIELD all-minimizers-proof
By \(\text{ass:coverage-level}\), \(c=1-\alpha>0\) and \(\tau=1+c=2-\alpha<2\).  The pair \((X,X)\) belongs to \(\mathcal F_{\tau}\), so the feasible-rectangle family is nonempty.  By \(\text{lem:hall-coverage}\), for every \(S\subseteq E\),
\[
\begin{aligned}
S\in\mathcal V_{\mu,c}
&\iff h_{\mu}(S)\geq c\\
&\iff \exists\,\varnothing\neq A,B\subseteq X:
A-B\subseteq S\ \text{ and }\ \mu_1(A)+\mu_0(B)-1\geq c\\
&\iff \exists\,(A,B)\in\mathcal F_{\tau}:A-B\subseteq S.
\end{aligned}
\]
The positive threshold \(c\) is used in the second equivalence to exclude the separate zero term in the Hall maximum.

Let
\[
r=\min_{(A,B)\in\mathcal F_{\tau}}|A-B|,
\]
which exists because the domain is a nonempty finite set.  Every feasible difference set \(A-B\) is valid by the displayed equivalence, so \(k^*\leq r\).  Conversely, if \(S\) is valid, the same equivalence gives a feasible difference set \(A-B\subseteq S\), whence \(r\leq|A-B|\leq|S|\).  Minimizing over valid \(S\) gives \(r\leq k^*\), and therefore
\[
k^*=r.
\]
If \(S\in\mathcal K^*\), choose a feasible \((A,B)\) with \(A-B\subseteq S\).  Then
\[
k^*\leq|A-B|\leq|S|=k^*,
\]
so \(A-B=S\).  Conversely, every feasible \(A-B\) of cardinality \(k^*\) is valid and belongs to \(\mathcal K^*\).  This proves the displayed characterization after identical subsets are deduplicated.

For inclusion minimality, suppose first that \(S\) is inclusion-minimal in \(\mathcal V_{\mu,c}\).  It contains a feasible difference set \(D=A-B\), which is itself valid; hence minimality forces \(S=D\).  Moreover, no distinct feasible difference set can be a proper subset of \(D\), since that subset would be valid.  Conversely, let \(D=A-B\) be an inclusion-minimal distinct feasible difference set.  If a proper subset \(T\subsetneq D\) were valid, then \(T\) would contain a feasible difference set \(D'\subseteq T\subsetneq D\), contradicting the choice of \(D\).  Thus \(D\) is inclusion-minimal valid.  All steps are finite and weak inequalities are retained, so zero cells, exact threshold equality, and nonunique certificates cause no exception.
%%% FIELD general-frontier-proof
Fix \(1\leq k\leq2L-1\).  By \(\text{thm:all-minimizers}\), a valid \(S\) with \(|S|\leq k\) exists if and only if there are nonempty \(A,B\subseteq X\) such that
\[
|A-B|\leq k
\quad\text{and}\quad
\mu_1(A)+\mu_0(B)\geq\tau.
\]
Because the maximum defining \(R_k\) is over a nonempty finite family, this is equivalent to \(R_k\geq\tau\).

For intervals, first suppose \(B_k\geq\tau\).  Finiteness gives lengths \(a,b\in\{1,\ldots,L\}\), with \(a+b\leq k+1\), and outcome intervals \(A,B\subseteq X\) of those lengths such that
\[
\mu_1(A)+\mu_0(B)=Q_1(a)+Q_0(b)\geq\tau.
\]
Their difference is the consecutive integer interval from \(\min A-\max B\) through \(\max A-\min B\), of cardinality \(a+b-1\leq k\).  It is valid by \(\text{thm:all-minimizers}\).

Conversely, suppose a consecutive effect interval \(I\subseteq E\) is valid and \(|I|\leq k\).  Choose a feasible rectangle \((A,B)\in\mathcal F_{\tau}\) with \(A-B\subseteq I\).  The sets are nonempty.  Let
\[
\overline A=\{\min A,\min A+1,\ldots,\max A\},
\qquad
\overline B=\{\min B,\min B+1,\ldots,\max B\}.
\]
The two endpoint differences \(\min A-\max B\) and \(\max A-\min B\) belong to \(A-B\subseteq I\).  Since \(I\) is consecutive,
\[
\overline A-\overline B
=\{\min A-\max B,\ldots,\max A-\min B\}
\subseteq I.
\]
Writing \(a=|\overline A|\) and \(b=|\overline B|\), this gives
\[
a+b-1=|\overline A-\overline B|\leq|I|\leq k.
\]
Because all marginal masses are nonnegative,
\[
Q_1(a)+Q_0(b)\geq\mu_1(\overline A)+\mu_0(\overline B)
\geq\mu_1(A)+\mu_0(B)\geq\tau,
\]
and hence \(B_k\geq\tau\).  This proves both profile equivalences.

The functions \(R_k\) and \(B_k\) are nondecreasing in \(k\).  Thus
\[
\exists k:\ B_k<\tau\leq R_k
\iff k^*<I^*,
\]
where the forward implication follows from \(k^*\leq k<I^*\), and the reverse implication uses \(k=k^*\).  A nonempty finite subset of consecutive integers is connected exactly when it is an interval.  If \(k^*<I^*\), no member of \(\mathcal K^*\) can be an interval.  Conversely, if every member of \(\mathcal K^*\) is disconnected but \(k^*=I^*\), an interval attaining \(I^*\) would itself belong to \(\mathcal K^*\), a contradiction.  Finally, the selector order first minimizes cardinality, so \(S^\star\in\mathcal K^*\), proving the last assertion.
%%% FIELD ternary-frontier-proof
Write \(p_i=\mu_{1,i}\), \(q_j=\mu_{0,j}\), \(m_1=\max_i p_i\), and \(m_0=\max_j q_j\).  We first give the exhaustive support classification.

For any finite nonempty integer sets \(A=\{a_1<\cdots<a_r\}\) and \(-B=\{b_1<\cdots<b_s\}\), the chain
\[
a_1+b_1<\cdots<a_r+b_1<a_r+b_2<\cdots<a_r+b_s
\]
contains \(r+s-1\) distinct points of \(A-B\).  Hence
\[
|A-B|\geq|A|+|B|-1. \tag{1}
\]
If \(|A-B|\leq2\), equation \((1)\) implies that, apart from singleton--singleton rectangles, exactly one factor is a singleton and the other has two points.  A two-point subset of \(X=\{0,1,2\}\) is either adjacent or is \(\{0,2\}\).  The largest Hall excess from adjacent two-point factors is
\[
\max\{m_1-\min(q_0,q_2),\ m_0-\min(p_0,p_2),0\}=b_2,
\]
because the largest mass of an adjacent pair under \(q\) is \(1-\min(q_0,q_2)\), and similarly under \(p\).  Singleton--singleton excesses are already dominated: for example
\[
m_1+m_0-1\leq m_1-\min(q_0,q_2),
\]
since \(\min(q_0,q_2)\leq1-m_0\); the symmetric inequality is identical.  The largest excess from a gapped factor is
\[
\max\{m_1-q_1,\ m_0-p_1,0\}=d_2. \tag{2}
\]
Indeed, the relevant rectangles are exactly
\[
\{i\}\times\{0,2\}
\quad\text{and}\quad
\{0,2\}\times\{j\},
\]
whose difference sets and excesses are, respectively,
\[
\{i-2,i\},\ p_i+q_0+q_2-1=p_i-q_1,
\]
and
\[
\{-j,2-j\},\ q_j+p_0+p_2-1=q_j-p_1. \tag{3}
\]

Next consider three-point difference sets.  Equation \((1)\) leaves factor-size patterns \((1,3)\), \((3,1)\), and \((2,2)\), besides patterns already producing at most two points.  The first two yield intervals.  If \(A\) and \(B\) both have two points with gaps \(g_A,g_B\in\{1,2\}\), their four formal differences have a collision exactly when \(g_A=g_B\).  Equal gap one gives a three-point interval; equal gap two forces \(A=B=\{0,2\}\) and gives the sole disconnected three-point difference set
\[
T=\{-2,0,2\}. \tag{4}
\]
Unequal gaps give four consecutive differences.  If one factor has all three points, pairing it with an adjacent two-point set gives four consecutive differences, pairing it with \(\{0,2\}\) gives all five differences, and pairing two three-point sets also gives all five.  Consequently every difference set of cardinality four or five is an interval.  This exhausts all nonempty pairs of subsets of \(X\), without a support-positivity assumption.

For later use, the interval profile can be evaluated explicitly.  Since
\[
Q_d(1)=m_d,\qquad Q_d(2)=1-\min(\mu_{d,0},\mu_{d,2}),\qquad Q_d(3)=1,
\]
the definition of \(B_3\) gives
\[
B_3-1
=\max\{m_1,m_0,1-\min(p_0,p_2)-\min(q_0,q_2)\}=b_3. \tag{5}
\]
All smaller interval certificates are dominated by these terms.  Also
\[
d_2\leq\max(m_1,m_0)\leq b_3. \tag{6}
\]

The relation corresponding to \(T\) connects the even block \(\{0,2\}\times\{0,2\}\) and the odd block \(\{1\}\times\{1\}\), with no cross-parity pair.  Monotonicity of marginal mass within each block and \(\text{lem:hall-coverage}\) therefore give
\[
h_{\mu}(T)=\max\{0,\ 1-p_1-q_1,\ p_1+q_1-1\}. \tag{7}
\]
If \(c>b_3\), then \(p_1<c\), \(q_1<c\), and, because \(c<1\),
\[
p_1+q_1-1<2c-1<c. \tag{8}
\]
Thus, in the regime \(c>b_3\), equation \((7)\) says that \(T\) is valid exactly when \(1-p_1-q_1\geq c\).

We can now classify all minima.  If \(d_2\geq c>b_2\), equation \((2)\) gives a valid disconnected two-point difference set, while the calculation of \(b_2\) excludes every interval of size at most two, including every singleton.  Hence \(k^*=2\), and equations \((2)\)--\((3)\) show that \(\mathcal K^*\) is exactly the two displayed, deduplicated families in the statement.  Conversely, if every cardinality minimizer is disconnected and \(k^*=2\), some gapped rectangle must have excess at least \(c\), while no adjacent rectangle may do so; hence \(d_2\geq c>b_2\).

If \(1-p_1-q_1\geq c>b_3\), equations \((5)\)--\((8)\) show that \(T\) is valid, no interval of size at most three is valid, and, by \((6)\), no set of size at most two is valid.  Thus \(k^*=3\) and the exhaustive classification forces \(\mathcal K^*=\{T\}\).  Conversely, if every cardinality minimizer is disconnected and \(k^*=3\), equation \((4)\) forces \(T\) to be the unique minimizer; absence of a valid interval of size at most three gives \(c>b_3\), and then \((7)\)--\((8)\) give \(1-p_1-q_1\geq c\).  Cardinality one cannot be disconnected, and the support enumeration shows that a cardinality-four or cardinality-five difference set is an interval, so there are no further cases.  Finally, \((6)\) proves that the two displayed regimes are mutually exclusive.  Every feasibility boundary was treated with \(\geq\), every exclusion with \(>\), so the conclusion includes zero cells and all threshold and selector ties.
%%% FIELD sharp-interval-price-proof
We first prove the deterministic additive-combinatorial inequality used below.  If \(A=\{a_1<\cdots<a_r\}\) and \(-B=\{b_1<\cdots<b_s\}\) are finite and nonempty, then
\[
a_1+b_1<\cdots<a_r+b_1<a_r+b_2<\cdots<a_r+b_s
\]
are \(r+s-1\) distinct members of \(A-B\).  Therefore
\[
|A-B|\geq|A|+|B|-1. \tag{1}
\]

If \(k^*=1\), the optimal singleton is itself an interval, so \(I^*=1\).  Suppose \(k^*=2\).  By \(\text{thm:all-minimizers}\), some \((A,B)\in\mathcal F_{\tau}\) has \(|A-B|=2\).  Equation \((1)\) gives \(|A|+|B|\leq3\).  The pattern \((1,1)\) would give only one difference, so the sizes are \((1,2)\) or \((2,1)\).  Thus \(A-B\) is a translate or reflected translate of a two-point subset of \(X\), and
\[
\operatorname{span}(A-B)\leq L-1.
\]
Its consecutive hull has at most \(L\) points and is valid because it contains the valid set \(A-B\).  Hence \(I^*\leq L\).  If \(k^*\geq3\), the full effect support \(E\) is valid and has \(2L-1\) points, so \(I^*\leq2L-1\).  These cases imply
\[
\frac{I^*}{k^*}\leq
\begin{cases}
1,&k^*=1,\\
L/2,&k^*=2,\\
(2L-1)/3,&k^*\geq3,
\end{cases}
\]
and
\[
I^*-k^*\leq
\begin{cases}
0,&k^*=1,\\
L-2,&k^*=2,\\
2L-4,&k^*\geq3.
\end{cases}
\]
Since \(L\geq3\), \(L/2\leq(2L-1)/3\) and \(L-2\leq2L-4\), proving both universal bounds.

For sharpness, fix \(0<\alpha<1/2\) and
\[
0<\varepsilon<\min\{\alpha/2,1-2\alpha\}.
\]
Let \(p=q\) have endpoint masses \(M=(1-\varepsilon)/2\) and each of its \(L-2\) interior masses equal to \(\varepsilon/(L-2)\).  The masses are strictly positive and sum to one.  With \(A=B=\{0,L-1\}\),
\[
A-B=\{-L+1,0,L-1\},
\qquad
p(A)+q(B)-1=1-2\varepsilon>1-\alpha=c.
\]
Thus this three-point set is valid.

Let \(S\) have at most two points.  It omits at least one member of \(\{-L+1,0,L-1\}\).  If it omits \(0\), the diagonal coupling \(\gamma_{ij}=p_i\mathbf 1\{i=j\}\) assigns zero mass to \(S\), so \(h_{\mu}(S)=0<c\).  If it omits \(-L+1\), put mass \(M\) at the corner \((0,L-1)\); if it omits \(L-1\), put mass \(M\) at \((L-1,0)\).  In either case, subtracting the corner mass leaves nonnegative row and column marginals of common total \(1-M\), which can be completed by their normalized product when \(1-M>0\).  The omitted corner then gives
\[
h_{\mu}(S)\leq1-M=\frac{1+\varepsilon}{2}<1-\alpha=c,
\]
where the strict inequality is exactly \(\varepsilon<1-2\alpha\).  Therefore \(k^*=3\).

Every proper consecutive subset of \(E\) omits at least one of the two extreme effects.  The corresponding corner construction gives the same strict upper bound \((1+\varepsilon)/2<c\).  The full support \(E\) has coverage one, so \(I^*=2L-1\).  Consequently
\[
\frac{I^*}{k^*}=\frac{2L-1}{3},
\qquad
I^*-k^*=2L-4.
\]
The construction proves attainment inside the stated strictly positive class; it does not identify any particular empirical trial law with the extremizer.
%%% FIELD exact-computation-proof
Fix \(S\subseteq E\), put \(p=\mu_1\) and \(q=\mu_0\), and construct a directed network with source \(s\), row vertices \(r_i\), column vertices \(c_j\), and sink \(t\).  Give \(s\to r_i\) capacity \(p_i\), give \(c_j\to t\) capacity \(q_j\), and include all \(L^2\) arcs \(r_i\to c_j\), with capacity \(2\) when \(i-j\notin S\) and capacity zero otherwise.  The network has \(2L+2\) vertices and exactly \(L^2+2L\) capacitated arcs.

The cut placing only \(s\) on the source side has capacity one.  Therefore no minimum cut crosses a row--column arc of capacity two.  If \(A\subseteq X\) is the set of source-side rows and \(C\subseteq X\) the set of source-side columns, a cut not crossing such an arc must satisfy
\[
N_{S^{\mathsf c}}(A)\subseteq C.
\]
For fixed \(A\), nonnegativity of \(q\) makes the smallest such cut take \(C=N_{S^{\mathsf c}}(A)\), and its capacity is
\[
\sum_{i\notin A}p_i+\sum_{j\in N_{S^{\mathsf c}}(A)}q_j
=1-p(A)+q(N_{S^{\mathsf c}}(A)).
\]
Hence the minimum-cut value \(\kappa_S\) satisfies, by \(\text{lem:hall-coverage}\),
\[
\kappa_S
=1-\max_{A\subseteq X}\{p(A)-q(N_{S^{\mathsf c}}(A))\}
=1-h_{\mu}(S).
\]
Thus \(\mathsf{Alg}\) returns \(h_{\mu}(S)=1-\kappa_S\) exactly.  This cut calculation remains valid with zero capacities and tied cuts.

There are \(2^L-1\) nonempty subsets of \(X\), so enumerating \(\mathcal F_{\tau}\) examines at most
\[
(2^L-1)^2
\]
pairs.  For each retained pair, the finite set \(A-B\) is computed, identical difference sets are deduplicated, and \(\text{thm:all-minimizers}\) proves that the retained sets of smallest cardinality are exactly \(\mathcal K^*\).  Applying the lexicographic order \(\prec\) then returns \(S^\star\).  Alternatively, direct enumeration has exactly \(2^{|E|}=2^{2L-1}\) candidate effect sets.  At \(L=7\), the two counts are
\[
(2^7-1)^2=127^2=16{,}129,
\qquad
2^{13}=8{,}192.
\]
All optimizations are finite; the word ``exact'' refers to exact comparison of the supplied marginal masses, not to a polynomial-time claim in \(L\).
%%% FIELD simultaneous-certification-proof
First record the partial-identification map.  Let the observed outcome be \(Y=ZY(1)+(1-Z)Y(0)\), and let \(P\) be the law of \((Y,Z)\) on \(X\times\{0,1\}\), a \(2L\)-cell probability law of dimension \(2L-1\).  By \(\text{ass:randomization}\) and \(\text{ass:assignment-overlap}\),
\[
P(Y=j\mid Z=d)=\Pr(Y(d)=j)=\mu_{d,j},
\qquad d\in\{0,1\},\ j\in X.
\]
Thus the two marginal vectors are identified.  The unobserved parameter is the \(L\times L\) joint-probability table \(\gamma\) of \((Y(1),Y(0))\).  With no cross-world restriction, its identified set is exactly
\[
\Theta_I(P)=\Gamma(\mu_1,\mu_0).
\]
Indeed, every compatible latent law has these margins, and conversely any table in \(\Gamma(\mu_1,\mu_0)\), combined with an independent treatment draw having the observed assignment probability, induces the same \(P\).  For fixed \(S\), the causal coverage is the linear functional \(\theta_S(\gamma)=\sum_{i-j\in S}\gamma_{ij}\); its identified set is a compact interval because \(\Gamma\) is compact and convex, and \(h_{\mu}(S)\) is its sharp lower endpoint, not the causal coverage under an arbitrarily chosen coupling.

We next derive the concentration inequality used here.  If \(W\) is Bernoulli with mean \(u\in[0,1]\), define
\[
f_u(\lambda)=\log(1-u+ue^\lambda)-u\lambda.
\]
Writing \(v=ue^\lambda/(1-u+ue^\lambda)\), direct differentiation gives
\[
f_u(0)=f_u'(0)=0,
\qquad
f_u''(\lambda)=v(1-v)\leq\frac14.
\]
Taylor's formula with integral remainder therefore yields \(f_u(\lambda)\leq\lambda^2/8\) for every real \(\lambda\), including the boundary cases \(u=0,1\).  For independent Bernoulli variables \(W_1,\ldots,W_n\) of mean \(u\), exponential Markov inequality and independence give, for \(t>0\) and \(\lambda>0\),
\[
\Pr\!\left\{\sum_{r=1}^n(W_r-u)\geq nt\right\}
\leq\exp\{-\lambda nt+n\lambda^2/8\}.
\]
Choosing \(\lambda=4t\), and applying the same argument to \(1-W_r\), gives
\[
\Pr\{|\overline W-u|>t\}\leq2e^{-2nt^2}. \tag{1}
\]

Under \(\text{ass:arm-sampling}\), each coordinate of \(\widehat\mu_d\) is an average of \(n_d\) independent Bernoulli indicators with mean \(\mu_{d,j}\).  At
\[
r_d=\sqrt{\frac{\log(4L/\beta)}{2n_d}},
\]
equation \((1)\) gives
\[
\Pr\{|\widehat\mu_{d,j}-\mu_{d,j}|>r_d\}
\leq2e^{-2n_dr_d^2}=\frac{\beta}{2L}.
\]
A union bound over both arms and all \(L\) coordinates proves, uniformly over the two simplices,
\[
\Pr\{\mu_0\in C_0,\ \mu_1\in C_1\}\geq1-\beta. \tag{2}
\]
The empirical vector belongs to \(C_d\), so each confidence polytope is nonempty; it is also a compact polytope.

For the optimization identity, let \(\mathcal G_S\) be the set of nonnegative tables \(\gamma\) whose row marginal \(r(\gamma)\) lies in \(C_1\) and whose column marginal \(s(\gamma)\) lies in \(C_0\).  Every \(\gamma\in\mathcal G_S\) belongs to \(\Gamma(r(\gamma),s(\gamma))\).  Conversely, for every \((\nu_1,\nu_0)\in C_1\times C_0\), every member of \(\Gamma(\nu_1,\nu_0)\) lies in \(\mathcal G_S\), and this coupling polytope is nonempty because it contains \(\nu_{1,i}\nu_{0,j}\).  Hence, without interchanging a minimum and a maximum,
\[
\begin{aligned}
\ell(S)
&=\min_{\nu_1\in C_1,\nu_0\in C_0}
\min_{\gamma\in\Gamma(\nu_1,\nu_0)}
\sum_{i-j\in S}\gamma_{ij}\\
&=\min_{\gamma\in\mathcal G_S}\sum_{i-j\in S}\gamma_{ij}.
\end{aligned} \tag{3}
\]
The final problem has linear objective and linear nonnegativity, marginal, simplex, and confidence-box constraints, and is therefore the asserted finite linear program.  Compactness gives attainment.

On the event in \((2)\), the true pair belongs to \(C_1\times C_0\), so equation \((3)\) gives simultaneously for every one of the finitely many \(S\subseteq E\),
\[
\ell(S)\leq h_{\mu}(S).
\]
Therefore \(S\in\widehat{\mathcal H}\) implies \(h_{\mu}(S)\geq\ell(S)\geq c\).  The implication holds for the entire random family on one common event, so choosing any data-dependent subfamily introduces no further multiplicity term.  Combining with \((2)\) proves the final uniform probability bound.
%%% FIELD uniform-selection-proof
We first establish a uniform modulus.  For any two marginal pairs \((p,q)\) and \((p',q')\), \(\text{lem:hall-coverage}\) and the elementary inequality \(|\max_a x_a-\max_a y_a|\leq\max_a|x_a-y_a|\) give, uniformly in \(S\),
\[
|h_{(q,p)}(S)-h_{(q',p')}(S)|
\leq \sup_{A\subseteq X}|p(A)-p'(A)|
+\sup_{C\subseteq X}|q(C)-q'(C)|.
\]
For probability vectors, the supremum over subsets is total variation, equal to one half of the \(\ell_1\) distance.  Consequently
\[
|h_{(q,p)}(S)-h_{(q',p')}(S)|
\leq\frac L2\bigl(\|p-p'\|_\infty+\|q-q'\|_\infty\bigr). \tag{1}
\]

Set \(t=\eta/(4L)\), and let \(\mathcal E_n\) be the event that both empirical marginal vectors are within \(t\) of their truths in sup norm.  For completeness, the needed tail bound follows directly as follows.  For a Bernoulli variable of mean \(u\),
\[
\log\mathbb E e^{\lambda(W-u)}
=\log(1-u+ue^\lambda)-u\lambda\leq\lambda^2/8,
\]
because the second derivative of the middle expression is \(v(1-v)\leq1/4\) and its value and first derivative vanish at zero.  Independence, exponential Markov inequality, and optimization at \(\lambda=4t\) give
\[
\Pr\{|\widehat\mu_{d,j}-\mu_{d,j}|>t\}\leq2e^{-2n_dt^2}.
\]
The union bound therefore yields
\[
\Pr(\mathcal E_n^{\mathsf c})
\leq4L e^{-2n_{\min}t^2}
=4L\exp\!\left\{-\frac{n_{\min}\eta^2}{8L^2}\right\}. \tag{2}
\]

Assume now that \(r_0,r_1\leq t\), as in the theorem, and work on \(\mathcal E_n\).  If \(\nu_d\in C_d\), the triangle inequality and the definition of \(C_d\) give
\[
\|\nu_d-\mu_d\|_\infty
\leq\|\nu_d-\widehat\mu_d\|_\infty
+\|\widehat\mu_d-\mu_d\|_\infty
\leq2t=\frac{\eta}{2L}. \tag{3}
\]
Combining \((1)\) and \((3)\), every pair in \(C_1\times C_0\) satisfies
\[
|h_{(\nu_0,\nu_1)}(S)-h_{\mu}(S)|\leq\frac\eta2
\quad\text{for every }S\subseteq E.
\]
Taking the minimum over the confidence polytopes preserves the two-sided band:
\[
|\ell(S)-h_{\mu}(S)|\leq\frac\eta2
\quad\text{simultaneously for every }S\subseteq E. \tag{4}
\]
Here nonemptiness is guaranteed because \(\widehat\mu_d\in C_d\).

For \(\mu\in\mathcal M_{L,\eta}\), \(\text{ass:selected-validity-gap}\) and \((4)\) imply
\[
\ell(S^\star)\geq h_{\mu}(S^\star)-\eta/2\geq c+\eta/2>c,
\]
so \(S^\star\) is certified.  For every \(T\prec S^\star\), \(\text{ass:predecessor-invalidity-gap}\) and \((4)\) imply
\[
\ell(T)\leq h_{\mu}(T)+\eta/2\leq c-\eta/2<c,
\]
so no predecessor is certified.  The lexicographic definition of \(\widehat S\) then gives \(\widehat S=S^\star\) throughout \(\mathcal E_n\).  Equation \((2)\) proves the stated probability bound uniformly over the class.  No active-cut uniqueness was used.

Finally, along the limit in \(\text{ass:allocation-limit}\), write \(N=n_0+n_1\to\infty\).  Since \(n_0/N\to\rho\in(0,1)\), both \(n_0\) and \(n_1\), hence \(n_{\min}\), diverge.  The right side of \((2)\) therefore tends to zero.  The fixed positive margin \(\eta\) is essential only for selector recovery; no such margin was used for finite-sample certification.
%%% FIELD equality-regimes-statement
Fix \(L\geq3\), \((\mu_0,\mu_1)\in\Delta_L^2\), and \(0<\alpha<1\).  The following conditions are equivalent:
\[
\frac{I^*}{k^*}=\frac{2L-1}{3};
\qquad
I^*-k^*=2L-4;
\qquad
k^*=3\ \text{ and }\ I^*=2L-1;
\]
\[
R_2<\tau\leq R_3\ \text{ and }\ B_{2L-2}<\tau;
\]
and
\[
R_2<\tau\leq R_3
\quad\text{and}\quad
\min\{\mu_{1,0},\mu_{1,L-1},\mu_{0,0},\mu_{0,L-1}\}>\alpha.
\]
Thus either sharp interval-price equality forces both the multiplicative and additive equalities simultaneously.  For fixed \(L\geq3\), strictly positive nonsymmetric marginal pairs satisfying equality exist if and only if \(0<\alpha<1/2\); hence simultaneous equality is possible outside the symmetric endpoint-heavy family.
%%% FIELD equality-regimes-proof
We begin with the equality cases in the universal proof.  If \(k^*=1\), then \(I^*=1\).  If \(k^*=2\), the difference-set argument in \(\text{thm:sharp-interval-price}\) gives \(I^*\leq L\), and for \(L\geq3\),
\[
\frac L2<\frac{2L-1}{3},
\qquad
L-2<2L-4. \tag{1}
\]
If \(k^*\geq3\), then \(I^*\leq2L-1\), so
\[
\frac{I^*}{k^*}\leq\frac{2L-1}{k^*}\leq\frac{2L-1}{3},
\qquad
I^*-k^*\leq2L-1-k^*\leq2L-4. \tag{2}
\]
Equality in either chain in \((2)\) holds exactly when \(k^*=3\) and \(I^*=2L-1\).  Together with \((1)\), this proves equivalence of the first three conditions and shows that one equality cannot occur without the other.

By \(\text{thm:general-frontier}\),
\[
k^*=3\iff R_2<\tau\leq R_3. \tag{3}
\]
The full interval \(E\) is always valid.  Again by the interval part of that theorem,
\[
I^*=2L-1\iff B_{2L-2}<\tau. \tag{4}
\]
The concentration profiles are nondecreasing in their width.  Under the constraint \(a+b\leq2L-1\), every pair \((a,b)\) is dominated either by \((L,L-1)\) or by \((L-1,L)\).  Since
\[
Q_d(L)=1,
\qquad
Q_d(L-1)=1-\min\{\mu_{d,0},\mu_{d,L-1}\},
\]
we obtain
\[
\begin{aligned}
B_{2L-2}
&=\max\{1+Q_0(L-1),Q_1(L-1)+1\}\\
&=2-\min\{\mu_{1,0},\mu_{1,L-1},\mu_{0,0},\mu_{0,L-1}\}.
\end{aligned} \tag{5}
\]
Because \(\tau=2-\alpha\), equations \((3)\)--\((5)\) prove the remaining equivalences.  This is a computable finite characterization: \(R_2\) and \(R_3\) are maxima over the explicitly specified subset pairs in \(\text{def:frontier-profiles}\).

Equation \((5)\) also settles existence under positivity.  Necessity follows because equality implies both endpoint masses of each arm exceed \(\alpha\); their sum is at most one, so \(2\alpha<1\).  For sufficiency, fix \(0<\alpha<1/2\), choose
\[
0<\varepsilon<\min\{\alpha/4,(1-2\alpha)/4\},
\qquad
M=\frac{1-\varepsilon}{2},
\qquad
0<\delta<M-\alpha,
\]
and define the control marginal \(q=\mu_0\) and treated marginal \(p=\mu_1\) by
\[
q_0=q_{L-1}=M,
\qquad
p_0=M+\delta,
\qquad
p_{L-1}=M-\delta,
\]
\[
p_i=q_i=\frac{\varepsilon}{L-2},
\qquad 1\leq i\leq L-2.
\]
They are normalized, strictly positive, and nonsymmetric because \(p\neq q\).  All four endpoint masses exceed \(\alpha\), so \((5)\) gives \(B_{2L-2}<\tau\).  The rectangle \(A=B=\{0,L-1\}\) has three-point difference set and mass sum
\[
p(A)+q(B)=2(1-\varepsilon)>2-\alpha=\tau,
\]
so \(R_3\geq\tau\).

It remains to verify \(R_2<\tau\), rather than assume it.  The interior mass \(\varepsilon/(L-2)\) is smaller than \(\alpha<M-\delta\), so the two largest atoms in either arm are its two endpoints and every two-point subset has mass at most \(1-\varepsilon\).  The integer inequality \(|A-B|\geq|A|+|B|-1\) shows that \(|A-B|\leq2\) implies \(|A|+|B|\leq3\).  Hence every rectangle entering \(R_2\) has mass sum at most
\[
(M+\delta)+(1-\varepsilon)
<(2M-\alpha)+(1-\varepsilon)
=2-2\varepsilon-\alpha
<2-\alpha=\tau.
\]
Thus \(R_2<\tau\).  Equations \((3)\)--\((5)\) now give \(k^*=3\), \(I^*=2L-1\), and both equalities.  This explicit asymmetric family proves possibility outside the symmetric endpoint-heavy construction for every admissible \(L\) and \(\alpha<1/2\).
